% SM Table: Validation Datasets Summary
% Generated for V8 comprehensive validation section
% Package refresh: 2026-08-12

\begin{table}[!htbp]
\centering
\caption{Summary of validation datasets used for docking protocol assessment. \textit{Multiple complementary validation strategies provide evidence for docking reliability: external benchmark (DEKOIS), positive-control enrichment (MMV), and pose reproduction (redocking).}}
\label{tab:sm_validation_datasets}
\small
\begin{tabular}{lcccp{5cm}}
\toprule
\textbf{Target} & \textbf{PDB} & \textbf{Actives} & \textbf{Decoys/Negatives} & \textbf{Validation Type} \\
\midrule
\multicolumn{5}{l}{\textit{External benchmark validation:}} \\
PfDHFR & 7F3Y & 40 & 1,200 & DEKOIS 2.0 (property-matched) \\
\midrule
\multicolumn{5}{l}{\textit{Positive-control enrichment:}} \\
PfDHFR & 7F3Y & 399 & 399 & MMV Malaria Box (consensus) \\
PfCRT & 6UKJ & 359 & 40 & MMV (score-stratified) \\
PfATP4 & 9N10 & 99 & 99 & MMV (score-stratified) \\
\midrule
\multicolumn{5}{l}{\textit{Redocking validation:}} \\
PfDHFR & 7F3Y & \multicolumn{2}{c}{2 ligands, 1 success (MTX failed: grid on NADPH)} & Full-ligand redocking \\
PfCRT & 6UKJ & \multicolumn{2}{c}{1 ligand (1 success)} & Full-ligand redocking \\
PfClpP & 2F6I & \multicolumn{2}{c}{1 ligand, peptide fragment (RMSD 1.65 \si{\angstrom})} & Fragment pose reproduction \\
PfATP4 & 9N10 & \multicolumn{2}{c}{1 ligand, ADP cofactor (RMSD 1.23 \si{\angstrom})} & Cofactor pose reproduction \\
\bottomrule
\multicolumn{5}{p{14cm}}{
\noindent\textbf{Notes:} DEKOIS 2.0 provides independent external validation with carefully designed property-matched decoys, addressing trivial separation artifacts. MMV Malaria Box positive-control uses experimentally confirmed antimalarials; PfCRT and PfATP4 labels are score-derived (top/bottom 25\%) and measure retrodictive consistency. Redocking assesses pose reproduction accuracy (RMSD < 2.0 \si{\angstrom}). PfClpP (2F6I) and PfATP4 (9N10) have no co-crystallized small-molecule inhibitor; their ``1 success'' reflects fragment/cofactor pose reproduction, not full-ligand redocking, and should not be interpreted as equivalent to the full-ligand successes for PfDHFR and PfCRT. The MTX redocking failure (RMSD $\approx$ \SI{30}{\angstrom}) may partly reflect grid-positioning bias (NADPH-adjacent site) rather than scoring-function limitations alone. These complementary approaches establish docking protocol reliability while transparently reporting limitations (DEKOIS ROC-AUC 0.450 baseline).}
\end{tabular}
\end{table}
